\section{The law of large numbers and the stability of averages}

Let
\[
X_1,X_2,\ldots,X_n
\]
be independent observations from the same distribution, with
\[
\E[X_i]=\mu
\qquad\text{and}\qquad
\Var(X_i)=\sigma^2<\infty.
\]
The sample mean is
\[
\overline X_n=\frac1n\sum_{i=1}^nX_i.
\]
The subscript emphasizes that the statistic changes when the sample size changes.

From the previous chapter,
\[
\E[\overline X_n]=\mu
\]
and
\[
\Var(\overline X_n)=\frac{\sigma^2}{n}.
\]
The centre of the sampling distribution stays at \(\mu\), while its variance
decreases as \(n\) increases.

This already suggests a fundamental regularity: averages from large samples
should usually lie close to the population mean.

\subsection*{Randomness remains, but averages become stable}

Increasing the sample size does not make the individual observations less
random. Each \(X_i\) still follows the same probability law. What changes is the
behaviour of their average.

For a small sample, one unusually large or small observation may strongly affect
\(\overline X_n\). In a large sample, the same observation is combined with many
others and has less influence on the final average.

The reduction in variance gives the scale of this stabilization:
\[
\operatorname{SD}(\overline X_n)=\frac{\sigma}{\sqrt n}.
\]
Thus doubling the sample size does not halve the standard deviation of the
sample mean. To halve it, the sample size must be multiplied by four.

\begin{examplebox}
Suppose individual observations have standard deviation \(20\). Then
\[
\operatorname{SD}(\overline X_{25})=\frac{20}{5}=4,
\]
while
\[
\operatorname{SD}(\overline X_{100})=\frac{20}{10}=2.
\]
The observations themselves remain equally variable, but averages based on one
hundred observations fluctuate less than averages based on twenty-five.
\end{examplebox}

\subsection*{The law of large numbers}

The law of large numbers makes the stabilization precise. For every fixed
\(\varepsilon>0\),
\[
P\left(|\overline X_n-\mu|>\varepsilon\right)
\longrightarrow 0
\qquad\text{as }n\to\infty.
\]
Equivalently,
\[
P\left(|\overline X_n-\mu|\leq\varepsilon\right)
\longrightarrow 1.
\]

\begin{theorembox}
For independent identically distributed observations with finite mean \(\mu\) and
finite variance \(\sigma^2\), the sample mean converges in probability to the
population mean:
\[
\overline X_n\xrightarrow{P}\mu.
\]
This statement is the finite-variance form of the \textbf{law of large numbers}.
More general versions require only a finite expectation, but their proof needs
methods beyond the variance argument used here.
\end{theorembox}

The notation
\[
\overline X_n\xrightarrow{P}\mu
\]
means that the probability of a noticeable difference between \(\overline X_n\)
and \(\mu\) becomes small when the sample size becomes large.

The theorem does not say that the sample mean eventually becomes equal to
\(\mu\). It says that substantial deviations become increasingly unlikely.

\subsection*{Why the result is plausible}

Chebyshev's inequality gives
\[
P\left(|\overline X_n-\mu|\geq\varepsilon\right)
\leq
\frac{\Var(\overline X_n)}{\varepsilon^2}.
\]
Using
\[
\Var(\overline X_n)=\frac{\sigma^2}{n},
\]
we obtain
\[
P\left(|\overline X_n-\mu|\geq\varepsilon\right)
\leq
\frac{\sigma^2}{n\varepsilon^2}.
\]
The right-hand side tends to zero as \(n\) increases.

This short argument shows the essential mechanism. The sampling distribution of
the mean becomes concentrated because its variance decreases like \(1/n\).

\begin{remarkbox}
The law of large numbers is not a statement that random errors cancel perfectly.
It is a probability statement saying that large deviations of the average become
less likely as the sample size grows.
\end{remarkbox}

\subsection*{Sample proportions as a special case}

For Bernoulli observations,
\[
X_i=
\begin{cases}
1,&\text{success},\\
0,&\text{failure},
\end{cases}
\qquad
P(X_i=1)=p,
\]
the sample proportion is
\[
\widehat p_n=\frac1n\sum_{i=1}^nX_i.
\]
Because \(\E[X_i]=p\), the law of large numbers gives
\[
\widehat p_n\xrightarrow{P}p.
\]

This explains why relative frequencies are used to estimate probabilities. In a
long sequence of independent repetitions, the observed proportion of successes
becomes increasingly stable around the underlying success probability.

For example, if a mechanism has success probability \(p=0.4\), small samples may
produce proportions such as
\[
0.2,\qquad 0.5,\qquad 0.7.
\]
With larger samples, proportions far from \(0.4\) become less common.

\subsection*{What the law does and does not provide}

The law of large numbers provides the first foundation for estimation. If
\[
\overline X_n\xrightarrow{P}\mu,
\]
then a large-sample mean is a reasonable estimator of \(\mu\). Similarly,
\(\widehat p_n\) is a reasonable estimator of \(p\).

But the law does not answer several important questions:

\begin{itemize}
    \item How large must \(n\) be in a particular problem?
    \item What is the approximate probability of a specified estimation error?
    \item What is the shape of the remaining random fluctuations?
    \item How should uncertainty be reported after one sample is observed?
\end{itemize}

The law tells us that the estimator becomes close to the parameter. It does not
yet describe the detailed pattern of the difference
\[
\overline X_n-\mu.
\]
That is the role of the central limit theorem.

\begin{summarybox}
The law of large numbers explains why averages and proportions can be used to
learn about a probability model. As the sample size grows, their sampling
distributions become concentrated around the corresponding population
parameters. This produces stability, but it does not yet quantify the full shape
of the remaining uncertainty.
\end{summarybox}
